\cleardoublepage % memastikan bab baru dimulai di halaman ganjil (kanan)
\chapter{PENDAHULUAN}
\label{chap:pendahuluan}
% --- Latar Belakang ---
\section{Latar Belakang}

Di era transformasi digital, pengembangan infrastruktur fasilitas publik menuntut standar efisiensi operasional dan keamanan yang semakin tinggi. Hal ini selaras dengan visi Gedung ITB Innovation Park (IIP) yang saat ini tengah bertransisi menuju konsep \textit{Smart and Green Building}. Transisi ini didukung oleh regulasi pemerintah pusat, khususnya Peraturan Menteri Pekerjaan Umum dan Perumahan Rakyat Nomor 10 Tahun 2023 tentang Bangunan Gedung Cerdas \autocite{kemenpupr2023}. Sebagai pusat inkubasi bisnis dan inovasi riset di bawah naungan Direktorat Kawasan Sains dan Teknologi (DKST), gedung IIP memiliki kompleksitas manajerial yang tinggi. Gedung ini menampung berbagai entitas pengguna, mulai dari karyawan tetap direktorat, staf dari puluhan perusahaan rintisan (\textit{tenant}), hingga tamu harian, dengan proyeksi kapasitas maksimal mencapai 1.200 orang per hari. 

Manajemen akses fisik menjadi salah satu elemen paling krusial dalam operasional gedung dengan kepadatan tinggi tersebut. Berdasarkan hasil observasi dan wawancara dengan pengelola IIP, pengelolaan akses saat ini masih menghadapi tantangan besar pada aspek interoperabilitas dan integrasi. Perangkat keras pengaman seperti gerbang pembatas (\textit{turnstile gate}) memang telah diimplementasikan di lobi utama dan area basemen, serta telah saling terkoneksi pada satu jaringan. Dalam operasionalnya, saat pendaftaran biometrik dilakukan, data tersebut dikirimkan ke sebuah komputer pusat yang menjalankan perangkat lunak bawaan vendor untuk kemudian didistribusikan ke gerbang lainnya. Namun, ekosistem perangkat lunak bawaan ini bersifat sangat tertutup (\textit{siloed}) dan kaku karena hanya berbasis \textit{desktop} lokal. Akibatnya, data pergerakan personel dan log riwayat akses terkunci di dalam aplikasi komputer tersebut dan tidak dapat ditarik secara \textit{real-time} oleh platform \textit{Building Management System} (BMS) eksternal yang sedang dikembangkan oleh gedung. Kondisi ini memicu inefisiensi operasional dan fleksibilitas, karena petugas keamanan harus selalu bergantung secara fisik pada satu stasiun komputer spesifik tersebut untuk sekadar memantau okupansi gedung maupun melakukan manajemen kredensial. Ketergantungan pada solusi lokal vendor inilah yang menghambat terwujudnya sistem kendali gedung cerdas yang holistik, terpusat, dan dapat diakses dari mana saja \autocite{buildings14113446}.

Ketiadaan sistem manajemen akses yang terpusat ini menimbulkan urgensi operasional dan risiko keamanan yang serius. Dari perspektif manajerial, pencatatan absensi atau pergerakan karyawan menjadi tidak akurat; sebagai contoh, ketiadaan sensor keluar di area basemen menyebabkan hilangnya data rekam jejak penghuni yang meninggalkan gedung. Dari perspektif keamanan dan mitigasi bencana, kelemahan ini berdampak fatal. Tanpa adanya log pergerakan (\textit{in/out}) yang diagregasi secara \textit{real-time}, petugas keamanan tidak memiliki visibilitas terhadap total okupansi atau jumlah orang yang masih berada di dalam gedung. Dalam skenario evakuasi darurat, seperti aktivasi \textit{Master Control Fire Alarm} (MCFA), ketiadaan data okupansi yang presisi akan sangat menghambat operasi penyelamatan \autocite{SAINI2022106}. Selain itu, ketiadaan rekam jejak pergerakan yang tersentralisasi membuat proses investigasi keamanan (misalnya dalam kasus pencarian barang hilang) memakan waktu yang sangat lama karena petugas harus meninjau rekaman CCTV secara manual.

Untuk menjawab permasalahan tersebut, modernisasi menggunakan teknologi \textit{Face Recognition} (pengenalan wajah) pada gerbang akses telah diusulkan dan mulai diimplementasikan oleh manajemen gedung. Berbeda dengan kartu RFID yang rentan dipindahtangankan, teknologi biometrik wajah menawarkan validitas identitas yang absolut dan mekanisme \textit{contactless} yang mencegah antrian fisik di lobi \autocite{wanjare2024}. Namun, implementasi perangkat keras cerdas ini memunculkan permasalahan baru pada lapisan perangkat lunak. Perangkat lunak bawaan pabrik umumnya bersifat kaku, tertutup, dan tidak dirancang untuk membagikan data secara \textit{real-time} ke platform eksternal atau \textit{dashboard} kustom yang sedang dikembangkan oleh tim integrasi BMS IIP. Dominasi solusi tertutup dari berbagai vendor ini menciptakan lanskap teknologi yang terfragmentasi dan memicu risiko \textit{vendor lock-in}, di mana pengelola gedung terjebak dalam satu ekosistem yang sulit untuk saling bertukar informasi secara terbuka \autocite{s19020229}. Keterbatasan ini menghambat terciptanya satu pusat kendali informasi gedung yang holistik.

Kelemahan pada aspek integrasi perangkat lunak inilah yang menjadi dasar pemikiran dalam tugas akhir ini. Penelitian ini merupakan bagian dari proyek kolaboratif pengembangan ekosistem BMS terintegrasi di Gedung IIP. Adapun ruang lingkup pada tugas akhir ini difokuskan secara spesifik pada perancangan arsitektur lapisan penengah (\textit{middleware}) dan antarmuka manajemen kontrol akses untuk entitas karyawan. Oleh karena itu, penelitian ini mengusulkan pengembangan Sistem Manajemen Karyawan Terpusat Berbasis Integrasi Perangkat Kontrol Akses pada Bangunan Gedung Cerdas yang mampu berkomunikasi langsung dengan memori terminal gerbang biometrik melalui protokol \textit{Intelligent Security API} (ISAPI). Melalui pendekatan ini, sistem mampu mendistribusikan data pendaftaran biometrik karyawan secara serentak ke seluruh gerbang dari satu antarmuka terpusat, serta menarik dan mengagregasi log pergerakan penghuni secara \textit{real-time}. Dengan demikian, kelemahan administratif dan urgensi keamanan dalam pengelolaan personel di Gedung IIP dapat diselesaikan secara komprehensif.

% --- Rumusan Masalah ---
\section{Rumusan Masalah}
Berdasarkan latar belakang yang telah diuraikan, rumusan masalah dalam tugas akhir ini adalah sebagai berikut:
\begin{enumerate}
    \item Bagaimana arsitektur sistem manajemen berbasis web dapat dirancang untuk berkomunikasi dengan perangkat gerbang biometrik secara langsung tanpa mengandalkan aplikasi desktop pihak ketiga?
    \item Bagaimana strategi sinkronisasi menggunakan protokol ISAPI untuk melakukan otomasi sinkronisasi data biometrik dan penarikan log pergerakan karyawan secara otomatis dari perangkat gerbang?
    \item Bagaimana mengembangkan antarmuka dashboard berbasis web yang memungkinkan pemantauan riwayat akses secara real-time sekaligus menyediakan mekanisme verifikasi visual otomatis untuk kejadian anomali dengan tetap menjaga privasi data karyawan?
\end{enumerate}

% --- Tujuan ---
\section{Tujuan}
Tujuan utama dari pelaksanaan tugas akhir ini adalah membangun Sistem Manajemen Karyawan Terpusat Berbasis Integrasi Perangkat Kontrol Akses pada Bangunan Gedung Cerdas untuk manajemen entitas karyawan di Gedung IIP guna mengatasi ketergantungan pada aplikasi lokal yang tertutup.
\begin{enumerate}
    \item Membangun basis data dan layanan backend (API) terpusat untuk mengonsolidasikan manajemen identitas dan kredensial karyawan.
    \item Mengembangkan modul middleware menggunakan protokol ISAPI sebagai jembatan integrasi untuk otomasi sinkronisasi data biometrik dan penarikan log pergerakan terminal gerbang secara real-time.
    \item Membangun aplikasi antarmuka web (dashboard) yang mengintegrasikan fitur pemantauan log pergerakan karyawan dan mekanisme verifikasi visual otomatis untuk kejadian anomali.
\end{enumerate}
Kriteria keberhasilan dari pelaksanaan tugas akhir ini adalah sebagai berikut:
\begin{enumerate}
    \item Pengelola dapat melakukan manajemen data profil karyawan (Tambah/Ubah/Hapus) secara terpusat melalui web, dan sistem berhasil menyinkronkan data tersebut ke perangkat gerbang biometrik melalui protokol ISAPI.
    \item Sistem berhasil mengekstraksi, mengagregasi, dan menyajikan data log pergerakan karyawan secara akurat dan real-time sesuai dengan riwayat transaksi fisik pada perangkat keras.
    \item Aplikasi dashboard dapat diakses melalui jaringan lokal gedung dan berhasil mengimplementasikan perlindungan privasi dengan menampilkan foto snapshot wajah secara eksklusif hanya pada log kejadian anomali (akses ditolak/tidak dikenal).
\end{enumerate}

% --- Batasan Masalah ---
\section{Batasan Masalah}
Ruang lingkup pelaksanaan tugas akhir ini dibatasi pada aspek perancangan dan implementasi sistem manajemen akses untuk entitas karyawan di Gedung IIP. Adapun batasan-batasan yang diambil dalam pelaksanaan tugas akhir ini adalah sebagai berikut:
\begin{enumerate}
    \item Sistem yang dikembangkan hanya akan terintegrasi dengan perangkat gerbang biometrik yang sudah ada di Gedung IIP, dan tidak mencakup integrasi dengan sistem keamanan lainnya seperti kamera CCTV atau sistem alarm kebakaran.
    \item Sistem yang dikembangkan akan menggunakan protokol ISAPI untuk komunikasi dengan perangkat gerbang biometrik, dan tidak akan mengembangkan protokol komunikasi baru atau melakukan modifikasi pada perangkat keras gerbang.
    \item Pengembangan antarmuka dashboard akan difokuskan pada fitur pemantauan log pergerakan karyawan dan mekanisme verifikasi visual otomatis, tanpa mengembangkan fitur tambahan seperti analisis perilaku atau prediksi keamanan.
    \item Tugas akhir ini dikerjakan secara berkelompok yang terdiri dari 3 orang mahasiswa, yaitu Axelius Davin dengan NIM 18222016, Muhammad Rifa dengan NIM 18222004, dan Natanael Steven dengan NIM 18222054. Penulis dalam hal ini berfokus pada pengembangan sistem manajemen akses karyawan.
    \item Pendaftaran karyawan pada sistem web saat ini terbatas pada sinkronisasi Nama, ID, dan Jenis Kelamin, serta tidak mencakup pengunggahan foto wajah dikarenakan keterbatasan teknis (multipart boundary) pada firmware perangkat yang digunakan (V4.38.0).
    \item Pelacakan riwayat pergerakan dan status kehadiran karyawan pada sistem ini sangat bergantung pada ketersediaan perangkat keras di lapangan. Sistem tidak mencakup pengolahan data dari terminal gerbang area \textit{basement} (karena keterbatasan perangkat keras yang hanya menggunakan tombol fisik atau \textit{push button} untuk arah keluar), serta terminal gerbang arah keluar di area lobi utama (dikarenakan belum didapatkannya hak akses kredensial untuk integrasi perangkat tersebut selama masa penelitian berlangsung). Oleh karena itu, ketidaklengkapan pencatatan log keluar (\textit{check-out}) merupakan batasan fisik dan administratif yang berada di luar ruang lingkup pengembangan perangkat lunak ini.
\end{enumerate}

% --- Metodologi Pengerjaan TA ---
\section{Metodologi}
Metodologi yang digunakan dalam pelaksanaan tugas akhir ini mengadopsi kerangka kerja \textit{Design Thinking} yang dipadukan dengan pendekatan Model Sistem (\textit{Prototyping Model}). \textit{Design Thinking} dipilih karena berpusat pada pemahaman mendalam terhadap masalah pengguna di lapangan, sementara Model Sistem memungkinkan pengembangan sistem integrasi antara perangkat lunak dan perangkat keras industrial secara iteratif. 

Proses ini dilaksanakan dalam siklus iterasi, dimulai dari pembuktian konsep (\textit{Proof of Concept}) hingga integrasi menyeluruh dan penyempurnaan performa. Adapun tahapan-tahapan pelaksanaannya disusun ke dalam 7 (tujuh) langkah sistematis sebagai berikut:

\begin{enumerate}
    \item Investigasi dan Pengumpulan Fakta (Tahap \textit{Empathize}): Tahap ini berfokus pada pemahaman masalah pengguna melalui observasi langsung di lobi Gedung ITB Innovation Park (IIP) dan wawancara dengan pengelola gedung. Tujuannya adalah memetakan alur masuk manual saat ini, mengidentifikasi hambatan dalam pengelolaan akses biometrik yang terfragmentasi, dan mengonfirmasi ketiadaan visibilitas data pergerakan secara terpusat.
    \item Studi Literatur: Melakukan pencarian, pengelompokan, dan penapisan referensi ilmiah serta regulasi yang relevan untuk membangun landasan teori pendukung. Kajian difokuskan pada standar Bangunan Gedung Cerdas (BGC) yang mengacu pada Permen PUPR Nomor 10 Tahun 2023, tantangan interoperabilitas \textit{Internet of Things} (IoT), serta protokol komunikasi \textit{Intelligent Security API} (ISAPI).
    \item Analisis Kebutuhan (Tahap \textit{Define}): Berdasarkan informasi yang disintesis dari tahap \textit{Empathize} dan studi literatur, dilakukan definisi masalah secara spesifik. Analisis mendalam dilakukan terhadap kebutuhan fungsional (seperti manajemen profil, sinkronisasi otomatis, dan penarikan log) serta kebutuhan non-fungsional (seperti keamanan jaringan lokal dan privasi data).
    \item Perancangan Sistem (Tahap \textit{Ideate}): Tahap perumusan gagasan dan perancangan arsitektur solusi teknis. Kegiatan meliputi pemodelan basis data menggunakan \textit{Entity Relationship Diagram} (ERD), perancangan logika \textit{middleware} ISAPI, serta desain antarmuka pengguna pada \textit{dashboard} menggunakan diagram UML. 
    \item Implementasi Sistem (Tahap \textit{Prototype}): Realisasi rancangan menjadi sistem fisik dan perangkat lunak yang terintegrasi. Tahap ini mencakup pembangunan layanan \textit{backend} menggunakan NestJS dan PostgreSQL, pengembangan modul integrasi langsung untuk mengeksekusi perintah perangkat keras biometrik, serta pembuatan antarmuka \textit{frontend} menggunakan kerangka kerja Next.js.
    \item Evaluasi Sistem (Tahap \textit{Test}): Pengujian terhadap Sistem yang telah dibangun untuk memvalidasi kinerjanya dalam lingkungan operasional nyata. Pengujian mencakup pengukuran akurasi sinkronisasi data biometrik, ketepatan logika, dan kecepatan visualisasi. 
    \item Penarikan Kesimpulan: Merangkum hasil pengembangan dan evaluasi berdasarkan kriteria keberhasilan yang telah ditetapkan pada awal penelitian. Tahap ini juga mencakup penyusunan saran dan rekomendasi pengembangan untuk masa mendatang.
\end{enumerate}

% --- Sistematika Penulisan ---
\section{Sistematika Penulisan}
Laporan tugas akhir ini disusun dengan sistematika penulisan sebagai berikut:
\begin{enumerate}
    \item	Bab I Pendahuluan: Berisi latar belakang masalah pengelolaan akses karyawan di Gedung IIP, rumusan masalah dalam bentuk pertanyaan penelitian, tujuan utama dan detail, kriteria keberhasilan, batasan masalah, metodologi pengembangan sistem menggunakan pendekatan Design Thinking, serta sistematika penulisan laporan.
    \item   Bab II Studi Literatur: Berisi landasan teori yang mendukung penelitian, meliputi standar Bangunan Gedung Cerdas mengacu pada Permen PUPR Nomor 10 Tahun 2023, konsep interoperabilitas IoT, arsitektur protokol ISAPI Hikvision, serta kerangka kerja Design Thinking sebagai metode pengembangan sistem.
    \item	Bab III Analisis: Berisi hasil investigasi dan analisis terhadap kondisi sistem kontrol akses saat ini di lobi Gedung IIP, identifikasi kebutuhan pengguna, serta spesifikasi kebutuhan fungsional dan non-fungsional sistem.
    \item Bab IV Perancangan Sistem Manajemen Karyawan Terpusat: Berisi rancangan sistem secara komprehensif yang terbagi ke dalam enam aspek utama, meliputi deskripsi umum sistem, rancangan arsitektur perangkat lunak dan jaringan, desain struktur basis data (ERD), perancangan logika program pada \textit{middleware}, desain antarmuka pengguna (\textit{UI/UX}), serta spesifikasi antarmuka komunikasi (RESTful API dan WebSocket).
    \item Bab V Implementasi Sistem Manajemen Karyawan Terpusat: Berisi detail realisasi rancangan menjadi sistem yang fungsional, meliputi pengaturan lingkungan infrastruktur, implementasi basis data, penulisan kode program \textit{middleware} menggunakan kerangka kerja NestJS, integrasi komunikasi perangkat keras melalui protokol ISAPI, pengembangan antarmuka pengguna (\textit{frontend}) menggunakan Next.js, serta pelapisan keamanan jaringan menggunakan NGINX.
    \item	Bab VI Evaluasi: Berisi prosedur pengujian fungsionalitas sistem menggunakan metode Black-box Testing, hasil pengukuran performa sinkronisasi data, serta evaluasi fitur verifikasi visual kejadian anomali untuk menjawab kriteria keberhasilan yang telah ditetapkan.            
    \item   Bab VII Kesimpulan dan Saran: Berisi kesimpulan akhir dari seluruh proses pengerjaan tugas akhir yang merujuk pada pencapaian tujuan, serta saran untuk pengembangan sistem di masa depan.
    \item   Daftar Pustaka: Berisi daftar referensi yang digunakan dalam penyusunan laporan tugas akhir.
    \item   Lampiran: Berisi dokumen-dokumen pendukung seperti kode program, data lengkap hasil pengujian, rekaman wawancara, dan lain-lain. 
\end{enumerate}